\section{Stokes公式}
\label{sec:mcalc-stokes}

一元微积分基本定理$\int_a^bf'(t)\,\mathrm dt=f(b)-f(a)$可以用前两节的语言重读：$f$是区间上的$0$形式，$f'\,\mathrm dt=df$是它的外微分，右端是$f$在边界$\{a,b\}$上带诱导定向符号$-1,+1$的“积分”。外微分和有向积分都已定义，现在可以问：对一般流形上的$(n-1)$形式$\omega$，$d\omega$在区域上的积分是否等于$\omega$在边界上的积分。本节先在固定条件下证明一般公式，再在欧氏区域中把它写成熟悉的三个经典公式。

\subsection{流形上的边界积分公式}
\label{sec:mcalc-stokes-general}

\begin{theorem}[Stokes公式]
\label{thm:mcalc-stokes}
设$M$是定向的$n$维带边界光滑流形，$\partial M$取诱导定向，$\iota:\partial M\to M$为包含映射。则对每个紧支撑的$(n-1)$形式$\omega$，
\begin{equation}
\int_Md\omega=\int_{\partial M}\iota^*\omega .
\label{eq:mcalc-stokes}
\end{equation}
若$\partial M=\varnothing$，右端为零。
\end{theorem}

\begin{proof}
两端都对$\omega$线性。取正向坐标覆盖与从属的单位分解$\{\psi_\alpha\}$，则$\omega=\sum_\alpha\psi_\alpha\omega$是有限和，所以只需对支撑落在单张正向坐标图内的$\omega$证明。这样的图要么是内部图，坐标在$\mathbb R^n$中取值；要么是边界图，坐标在$\mathbb H^n=\{x^1\leq0\}$中取值。在坐标中写
\[
\omega=\sum_{i=1}^n\omega_i\,dx^1\wedge\cdots\wedge\widehat{dx^i}\wedge\cdots\wedge dx^n ,
\]
记号$\widehat{dx^i}$表示略去这一项，$\omega_i$光滑且有紧支撑。由式\eqref{eq:mcalc-exterior-derivative}，$d\omega_i\wedge(\cdots)$中只有$\partial_i\omega_i\,dx^i$一项存活，把$dx^i$移到第$i$个位置需越过$i-1$个因子，所以
\[
d\omega=\sum_{i=1}^n(-1)^{i-1}\,\partial_i\omega_i\;dx^1\wedge\cdots\wedge dx^n .
\]

内部图的情形，$\int d\omega=\sum_i(-1)^{i-1}\int_{\mathbb R^n}\partial_i\omega_i\,\mathrm dx$。对第$i$项先沿$x^i$积分：$\int_{-\infty}^{\infty}\partial_i\omega_i\,\mathrm dx^i$是$\omega_i$在远端两处的值之差，紧支撑使其为零。所以左端为零。这样的图中没有边界点，右端也为零。

边界图的情形，积分区域是$\{x^1\leq0\}$。对$i\geq2$的项，沿$x^i$的积分仍在整条直线上进行，同理为零。对$i=1$的项，沿$x^1$在$(-\infty,0]$上积分，由微积分基本定理得$\omega_1(0,x^2,\ldots,x^n)$。所以
\[
\int_Md\omega=\int_{\mathbb R^{n-1}}\omega_1(0,x^2,\ldots,x^n)\,\mathrm dx^2\cdots\mathrm dx^n .
\]
另一方面，边界上$x^1\equiv0$，故$\iota^*dx^1=0$；$\omega$中凡含$dx^1$的项拉回后为零，只剩$\iota^*\omega=\omega_1(0,x^2,\ldots,x^n)\,dx^2\wedge\cdots\wedge dx^n$。由上节，$\partial\mathbb H^n$的诱导定向就是坐标$(x^2,\ldots,x^n)$的标准定向，所以右端恰为上式。$n=1$时边界是一点，同样的计算就是微积分基本定理连同符号约定。
\end{proof}

证明中的每个条件都有明确的用途。单位分解把整体问题化为单张坐标图内的问题，这一步需要求和有限，由紧支撑保证。在坐标图内，公式归结为对一个变量用微积分基本定理，远端的值为零依赖紧支撑，边界值的符号依赖诱导定向的“外法向在前”约定。

\paragraph{条件不可随意去掉}
去掉紧支撑，公式可能失败。取$M=[0,\infty)\subseteq\mathbb R$，通常定向，$\omega=f$为常函数$1$。$d\omega=0$，左端为零；边界$\{0\}$处外向量指向负方向，诱导符号为$-1$，右端为$-f(0)=-1$。两端不等，原因是$f$在无穷远处不衰减：区间“右端的边界项”跑到了无穷远。若定向反转，两端同时变号，所以诱导定向必须与$M$的定向配套使用。

\paragraph{两个推论}
无边界的紧流形$M$上，任何$(n-1)$形式的外微分积分为零：$\int_Md\eta=0$。由此，例~\ref{ex:mcalc-form-integral-examples}中$\int_{\mathbb S^1}\omega=2\pi\neq0$直接说明穿孔平面上的闭形式$\omega=(x\,dy-y\,dx)/(x^2+y^2)$不是恰当形式：否则把$\omega=df$拉回单位圆，积分应为零。

另一个推论是面积可以由边界计算。在平面单位圆盘上取$\omega=x\,dy$，$d\omega=dx\wedge dy$，左端为面积$\pi$。右端用$c(t)=(\cos t,\sin t)$计算，$c^*(x\,dy)=\cos^2t\,dt$，积分$\int_0^{2\pi}\cos^2t\,\mathrm dt=\pi$，两端相等。

\subsection{Green、Gauss与经典Stokes公式}
\label{sec:mcalc-classical-formulas}

欧氏空间中的三个经典公式，都是式\eqref{eq:mcalc-stokes}在选取适当的流形与形式次数后的特例。下面逐一核对区域、边界、形式次数与方向约定。这里假设所涉区域的闭包紧、边界光滑，于是它是带边界的紧流形，任何光滑形式限制在其上都紧支撑。

\paragraph{Green公式}
取平面区域$D$，$n=2$，标准定向；边界$\partial D$是若干条闭曲线，诱导定向使区域位于前进方向的左侧（外法向在前、切向在后构成正向基）。取$1$形式$\omega=P\,dx+Q\,dy$，$d\omega=(\partial_xQ-\partial_yP)\,dx\wedge dy$。式\eqref{eq:mcalc-stokes}成为
\[
\iint_D\bigl(\partial_xQ-\partial_yP\bigr)\,\mathrm dx\,\mathrm dy=\oint_{\partial D}P\,\mathrm dx+Q\,\mathrm dy ,
\]
即\term{Green公式}（Green's theorem）。左端是区域积分，因为标准定向下$dx\wedge dy$的积分就是普通的面积积分。

\paragraph{Gauss公式}
取$\mathbb R^3$中的区域$V$，$n=3$，标准定向。对向量场$F=(F^1,F^2,F^3)$取$2$形式
\[
\omega_F=F^1\,dy\wedge dz+F^2\,dz\wedge dx+F^3\,dx\wedge dy ,
\]
它在两个向量$u,v$上的值是$\det[F,u,v]=\langle F,u\times v\rangle$。由式\eqref{eq:mcalc-exterior-derivative}，$d\omega_F=(\partial_xF^1+\partial_yF^2+\partial_zF^3)\,dx\wedge dy\wedge dz$，系数是$F$的散度。在边界点取外单位法向$n$，以及使$(n,u,v)$为正向标准正交基的切向量$u,v$，则$u\times v=n$，故$\iota^*\omega_F(u,v)=\langle F,n\rangle$。这说明边界上的$\iota^*\omega_F$等于“$\langle F,n\rangle$乘面积元”。式\eqref{eq:mcalc-stokes}成为
\[
\iiint_V\operatorname{div}F\,\mathrm dV=\iint_{\partial V}\langle F,n\rangle\,\mathrm dA ,
\]
即Gauss公式，也称\term{散度定理}（divergence theorem）。这里的“面积元”暂按欧氏曲面面积理解，它与黎曼体积的关系在第~\ref{sec:mcalc-volume}节说明，散度的几何定义在第~\ref{sec:mcalc-operators}节给出。

\paragraph{经典Stokes公式}
取$\mathbb R^3$中的定向曲面$S$，$n=2$，带边界曲线$\partial S$。曲面的定向由单位法向场$n$给出：切向量$(u,v)$正向当且仅当$(n,u,v)$在$\mathbb R^3$中正向。取$1$形式$\omega=F^1dx+F^2dy+F^3dz$，计算得$d\omega=\omega_{\operatorname{curl}F}$，即以旋度为系数的上述$2$形式。式\eqref{eq:mcalc-stokes}成为
\[
\iint_S\langle\operatorname{curl}F,n\rangle\,\mathrm dA=\oint_{\partial S}F^1\,\mathrm dx+F^2\,\mathrm dy+F^3\,\mathrm dz .
\]
诱导定向的约定在这里化为右手规则：右手四指沿边界的正向弯曲时，拇指指向$n$。理由是在边界点，外法向$N$切于曲面、指向曲面外侧，边界切向$t$须使$(N,t)$是曲面的正向基，即$(n,N,t)$为$\mathbb R^3$的正向基；$n\times N=t$正是右手规则描述的方向。

三个公式的条件各不相同。Green公式中$D$可以有洞，边界的每条曲线都按“区域在左侧”定向，内边界因此是顺时针的。Gauss公式要求法向朝外。经典Stokes公式中曲面可以弯曲，只要能连续地选取法向；在M\"obius带这样不可定向的曲面上，公式没有意义。这些公式在欧氏坐标中的直接证明与应用可参考本章开头所引的卷二章节。
